%%%%%%%%%%%%%%%%%%%%%%%%%%%%%%%%%%%%%%%%%%%%%%%%%%%%%%%%%%%%%%%%%%%%%%
% Lớp 12 | Ứng dụng đạo hàm để khảo sát hàm số | Cực trị của hàm số
% Dạng 2D1?2-6: Tìm 𝑚 để hàm số, đồ thị hàm số các hàm số khác có cực trị thỏa mãn điều
% Số câu: 3  (sinh tự động bởi tach_cau.py; từ nay thầy sửa trực tiếp file này)
%%%%%%%%%%%%%%%%%%%%%%%%%%%%%%%%%%%%%%%%%%%%%%%%%%%%%%%%%%%%%%%%%%%%%%

%%%%%%%%%%%%%%%% ĐÚNG/SAI %%%%%%%%%%%%%%%%

% nguồn: Phan1_Lop_12_Dai_so_Chuong_1
\begin{ex}%[2D1V2-6]%[ma:13f81f83]
	Cho hàm số $y=\dfrac{x^2 + 2x + 3}{2x +  1}$
	\choiceTF 
		{Tập xác định của hàm số là  $\mathbb{R}$}
		{Phương  trình $y' = 0$ có hai nghiệm nguyên dương phân biệt}
		{\True Hàm số đã cho có  hai điểm cực trị và hai điểm cực trị này nằm về hai phía của trục tung}
		{\True Khi đồ thị hàm số có hai điểm cực trị thì đường thẳng đi qua hai điểm cực trị của đồ thị hàm số có phương trình $y= x  + 1$}
	 \loigiai{
		\begin{itemchoice}
			\itemch Tập xác định  $D = \mathbb{R}\setminus \left\{-\dfrac{1}{2}\right\} $.
			\itemch $y' = \dfrac{2x^2 + 2x - 4}{(2x + 1)^2}, y'= 0 \Leftrightarrow 2x^2 + 2x -  4= 0 \Leftrightarrow \hoac{&x = 1, y(1) = 2\\&x = -2, y(-2) = -1}$
			\itemch Đồ thị hàm số có hai điểm cực trị là $M(1;2)$ và $N(-2;-1)$ và hai điểm cực trị này  nằm về hai phía của trục tung
			\itemch Phương  trình đường thẳng đi qua hai điểm cực trị $M, N$ của đồ thị hàm số đã cho là: $y = x + 1$.\\
			\textbf{Chú  ý:}\\
			Áp dụng tính chất: Nếu $x_0$ là điểm cực trị của hàm số hữu tỷ $y = \dfrac{u(x)}{v(x)}$ thì giá trị cực trị tương ứng của hà số là $y_0=\dfrac{u(x_0)}{v(x_0)}= \dfrac{u'(x_0)}{v'(x_0)}$. Suy ra với bài toán trên ta có phương trình đường thẳng qua hai điểm cực trị của đồ  thị hàm số là $y' = \dfrac{(x^2 + 2x + 3)'}{(2x + 1)'}= x + 1$.
		\end{itemchoice}
	 }
\end{ex}

% nguồn: Phan1_Lop_12_Dai_so_Chuong_1
\begin{ex}%[2D1V2-6]%[ma:5747e636]
	Cho  hàm số $y = \dfrac{x^2 - 2mx + m + 2}{x -m}$, với $m$ là tham số
	\choiceTF 
		{\True Tập xác định của  hàm số là $\mathbb{R}$}
		{\True Có hai giá trị nguyên của tham  số $m$ để hàm số có hai điểm cực trị}
		{\True Hàm số đạt cực đại tại $x = -1$ khi  $m =\dfrac{1}{2}$}
		{\True Khi đồ thị hàm số có hai điểm cực trị thì đường thẳng đi qua hai điểm cực trị cử đồ thị hàm số  có phương trình $y = 2x -2m$}
	\loigiai{ 
		\begin{itemchoice}
			\itemch Hàm  số xác định khi $x - m \ne 0 \Leftrightarrow x \ne m$ nên tập xác định là $D= \mathbb{R} \setminus\{m\}$
			\itemch Đạo hàm $y'=\dfrac{x^2 - 2mx + 2m^2 - m -  2}{(x - m)^2}$. Để hàm số có hai điểm cực trị thì $y' = 0$ có hai nghiệm phân biệt khác $m$ hay $g(x) = x^2 - 2mx -2m^2 - m - 2$ có hai nghiệm phân biệt khác $m$.\\
			$\heva{&\Delta' > 0\\&g(m) \ne 0} \Leftrightarrow \heva{&-m^2 + m + 2 > 0\\&m^2 - m - 2 \ne 0}\Leftrightarrow  m \in (-1;2)$\\
			Vì $m$ nguyên nên $m = \{0;1\}$ nên có hai giá trị  nguyên của tham số $m$ thỏa mãn.
			\itemch Hàm  số đạt cực trị tại $x = -1$ thì $y'(-1) = 0 \Leftrightarrow 2m^2 + m - 1= 0 \Leftrightarrow \hoac{&m = -1\\&m = \dfrac{1}{2}}$\\
			Thử lại với $m = \dfrac{1}{2}$ thì  $y'=\dfrac{x^2 - x - 2}{x - \dfrac{1}{2}}$ và có bảng biến thiên như sau:\\
			Vậy với $m = \dfrac{1}{2}$  thỏa mãn yêu cầu bài toán.
			\itemch Cho hàm số $y= \dfrac{u(x)}{v(x)}$. Nếu hàm số có hai điểm cực trị thì phương  trình đường thẳng đi qua hai điểm cực trị có dạng $y = \dfrac{u'(x)}{v'(x)}$.\\
			Áp dụng vào bài toán ta  được $y= \dfrac{(x^2 - 2mx + m + 2)^2}{(x - m)^2}=2x -2m$
		\end{itemchoice}
	} 
\end{ex}


%%%%%%%%%%%%%%%% TRẢ LỜI NGẮN %%%%%%%%%%%%%%%%

% nguồn: Phan2_Lop_12_Dai_so_Chuong_1
\begin{bt}%[2D1V2-6]%[ma:beb8f16d]
	Cho hàm số $y=\dfrac{x^2+mx+m^2-2m-4}{x-2}\quad (1)$. Tìm $m$ để đồ thị hàm số đã cho có hai điểm cực  trị và hai điểm cực trị cách đều đường thẳng $\Delta\colon 2x+y+1=0$.
	\shortans{-9} 
	 \loigiai{
		Đạo hàm $y'=\dfrac{x^2-4x+4-m^2}{(x-2)^2}$. Dấu của  $y'$ là dấu của $g(x)=x^2-4x+4-m^2$.\\
		Hàm số có hai điểm cực trị $\Leftrightarrow$ phương trình $g(x)=0$ có hai nghiệm phân biệt khác $2$, điều  này tương đương
		$\heva{&\Delta'=4-4+m^2=m^2>0 \\&4-8+4-m^2\ne0}\Leftrightarrow  m\neq 0$.\\
		Nghiệm của phương  trình $g(x)=0$ là $x_1=2-m$ hoặc $x_2=2+m$.\\
		Suy ra hai điểm cực trị của đồ thị hàm số là  $A(2-m; 4-m)$, $B(2+m; 4+3m)$\\
		$\mathrm{d}(A,\Delta)=\dfrac{|9-3m|}{\sqrt{5}}$; $\mathrm{d}(B,\Delta)=\dfrac{|9+5m|}{\sqrt{5}}$.\\ 
		$\mathrm{d}(A,\Delta)=\mathrm{d}(B,\Delta)\Leftrightarrow |9-3m|=|9+5m|$ 
		 $\Leftrightarrow
		 \hoac{&9-3m=9+5m \\&9-3m=-9-5m}\Leftrightarrow
		\hoac{&m = 0 \\&m = -9}$\\ 
		 So với điều kiện $m\ne 0$ nhận $m=-9$.
	 }
\end{bt}
